\documentclass[10pt]{article}
\usepackage[T1]{fontenc}
\usepackage[utf8]{inputenc}
\usepackage[french]{babel}
\usepackage{amsmath,amssymb,mathrsfs,mathpazo}
\usepackage{geometry}
\geometry{a4paper,margin=2cm}
\usepackage{graphicx,array,tabularx,tabularray,colortbl,multirow,multicol,booktabs,enumitem,pifont,ulem,fancybox,diagbox,cancel}
\usepackage{pstricks,pstricks-add,pst-plot,pst-tree,pst-node,pst-text,pst-func,pst-math,pst-bspline,pst-3dplot,pst-eucl,pst-solides3d}
\let\pscalculate\undefined
\usepackage{pst-bezier}
\usepackage{tikz,pgfplots,tkz-tab}
\usetikzlibrary{arrows,arrows.meta,patterns,calc,positioning,trees,decorations.markings,shapes.geometric,shapes.misc,shapes.arrows,intersections,decorations.pathmorphing,decorations.pathreplacing,backgrounds,shadows,babel}
\pgfplotsset{compat=1.18}
\usepgfplotslibrary{fillbetween}
\usepackage[np]{numprint}
\usepackage{listings,makecell,pas-tableur,twoopt}
\let\barre\undefined
\usepackage[pstricks]{tablvar}
\definecolor{darkgreen}{RGB}{0,100,0}
\providecommand{\red}{\color{red}}
\providecommand{\blue}{\color{blue}}
\providecommand{\green}{\color{green}}
\providecommand{\black}{\color{black}}
\providecommand{\white}{\color{white}}
\providecommand{\gray}{\color{gray}}
\providecommand{\cyan}{\color{cyan}}
\providecommand{\magenta}{\color{magenta}}
\providecommand{\yellow}{\color{yellow}}
\providecommand{\violet}{\color{violet}}
\providecommand{\orange}{\color{orange}}
\providecommand{\euro}{\texteuro}
\providecommand{\EUR}{\texteuro}
\providecommand{\ieme}{\textsuperscript{e}}
\providecommand{\degres}{\ensuremath{^\circ}}
\providecommand{\degree}{\ensuremath{^\circ}}
\providecommand{\vv}[1]{\overrightarrow{#1}}
\providecommand{\Coord}[1]{#1}
\providecommand{\decosix}{}
\providecommand{\starredbullet}{\ensuremath{\bullet}}
\pagestyle{empty}
\setlength{\parindent}{0pt}
\usepackage{adjustbox}\DeclareUnicodeCharacter{2236}{:}\DeclareUnicodeCharacter{2190}{\ensuremath{\leftarrow}}
\providecommand{\R}{\mathbb{R}}
\providecommand{\N}{\mathbb{N}}
\providecommand{\Z}{\mathbb{Z}}
\providecommand{\Q}{\mathbb{Q}}
\providecommand{\D}{\mathbb{D}}
\providecommand{\C}{\mathbb{C}}
\providecommand{\e}{{\mathrm{e}}}
\providecommand{\vect}[1]{\overrightarrow{#1}}
\providecommand{\vectt}[1]{\overrightarrow{\mathrm{#1}}}
\providecommand{\vvt}[1]{\overrightarrow{\mathrm{#1}}}
\providecommand{\np}[2][]{{#2~#1}}
\providecommand{\vv}[1]{{\overrightarrow{#1}}}
\renewcommand{\d}{{\mathrm{d}}}
\providecommand{\degres}{{^\circ}}
\providecommand{\degre}{{^\circ}}
\providecommand{\blue}{\color{blue}}
\providecommand{\red}{\color{red}}
\providecommand{\enskip}{\quad}
\providecommand{\llbracket}{\lbrack\!\lbrack}
\providecommand{\rrbracket}{\rbrack\!\rbrack}
\providecommand{\euro}{{\text{€}}}
\providecommand{\no}{{\text{n°}}}
\providecommand{\interval}[2]{{[#1;#2]}}
\providecommand{\intervFO}[2]{{[#1;#2[}}
\providecommand{\dotfill}{\ldots}

\providecommand{\ds}{\displaystyle}
\providecommand{\textding}[1]{#1}
\providecommand{\Oij}{\(\mathrm{O};\vec{\imath},\vec{\jmath}\)}
\providecommand{\Oijk}{\(\mathrm{O};\vec{\imath},\vec{\jmath},\vec{k}\)}
\providecommand{\Ouv}{\(\mathrm{O};\vec{u},\vec{v}\)}
\providecommand{\intervFF}[2]{\ensuremath{[#1;#2]}}
\renewcommand{\euro}{\ifmmode\text{\texteuro}\else\texteuro\fi}
\DeclareUnicodeCharacter{00B0}{\ifmmode^{\circ}\else\textdegree\fi}
\DeclareUnicodeCharacter{20AC}{\ifmmode\text{\texteuro}\else\texteuro\fi}
% Énoncés : Éducation nationale. Transcription et corrigés : APMEP, Denis Vergès et contributeurs. Découpage et classement : CoopMaths, Sébastien Lozano, Éric Elter et contributeurs. Adaptation de la présentation : Maths Explorer.
% CC BY-SA : https://coopmaths.fr/www/about/
% https://coopmaths.fr/alea/static/e3c/2020/tex/e3c_2020_00_sujet55_1.tex
\begin{document}
\noindent\textit{APMEP / CoopMaths. CC BY-SA. Présentation adaptée par Maths Explorer.}\par\medskip
\medskip

\emph{Cet exercice est un QCM et comprend cinq questions. Pour chacune des questions, une seule
des quatre réponses proposées est correcte. Les questions sont indépendantes.}

\emph{Pour chaque question, indiquer le numéro de la question et recopier sur la copie la lettre
correspondante à la réponse choisie.}

\emph{Aucune justification n'est demandée mais il peut être nécessaire d'effectuer des recherches
au brouillon pour aider à déterminer votre réponse.}

\emph{Chaque réponse correcte rapporte un point. Une réponse incorrecte ou une question sans
réponse n'apporte ni ne retire de point.}

\medskip

\textbf{Question 1}

\medskip

Une équation de la tangente à la courbe représentative de la fonction exponentielle au point
d'abscisse 0 est :

\begin{tabularx}{\linewidth}{*{4}{X}}
\textbf{a.~~}$y=x+1$  &\textbf{b.~~}$y=\e x$ &\textbf{c.~~$y=\e^x $}& \textbf{d.~~}$y= x-1 $ .
\end{tabularx}

\smallskip

\textbf{Question 2}

\medskip

La fonction $f$ définie sur $\R$ par : $f(x) = \e^{-2x+6}$ admet pour dérivée la fonction $f'$ définie
sur $\R$ par :
\begin{tabularx}{\linewidth}{*{2}{X}}
\textbf{a.~~}$f'(x) = \e^{-2x+6}$  &\textbf{b.~~} $f'(x) = -2\e^{-2x+6}$\\\textbf{c.~~}$f'(x) = -2x\e^{-2x+6}$& \textbf{d.~~}$f'(x) = (-2x + 6)\e^{-2x+6}$ .
\end{tabularx}

\medskip

\textbf{Question 3}

\medskip
Dans le repère orthonormé \Oij, le vecteur $\vv{AB}$ représenté ci-dessous est égal à :

\begin{center}
\psset{labelFontSize=\scriptstyle,showorigin=false}
\begin{pspicture}(-0.75,-0.5)(7.5,3.5)
\multido{\n=0+1}{8}{\psline[linewidth=0.75pt,linecolor=lightgray](\n,0)(\n,3.2)}
\multido{\n=0+0.2}{37}{\psline[linewidth=0.5pt,linecolor=lightgray](\n,0)(\n,3.2)}
\multido{\n=0+1}{4}{\psline[linewidth=0.75pt,linecolor=lightgray](0,\n)(7.2,\n)}
\multido{\n=0+0.2}{17}{\psline[linewidth=0.5pt,linecolor=lightgray](0,\n)(7.2,\n)}
\psaxes[linewidth=0.95pt]{->}(0,0)(7.4,3.3)
\psdots[dotstyle=+,dotscale =1.4,dotangle=45](1,3)(7,1)
\uput[l](1,3){A} \uput[dl](7,1){B}\psset{arrowscale=1.8}
\psline[linewidth=0.8pt,linecolor=darkgray]{->}(0,0)(0,1)\psline[linewidth=0.8pt,linecolor=darkgray]{->}(0,0)(1,0)
\psline[linewidth=0.8pt,linecolor=darkgray]{->}(1,3)(7,1)
\uput[dl](0,0){O}\uput[d](0.5,0){$\vv{\imath}$}\uput[l](0,0.5){$\vv{\jmath}$}

\end{pspicture}
\end{center}

\medskip

\begin{tabularx}{\linewidth}{*{4}{X}}
\textbf{a.~~}$-2\vv{\imath} + 6 \vv{\jmath}$  &\textbf{b.~~}$-6\vv{\imath}+ 2\vv{\jmath}$ &\textbf{c.~~}$2\vv{\imath} -6\vv{\jmath}$& \textbf{d.~~}$6\vv{\imath}- 2\vv{\jmath}$ .
\end{tabularx}

\medskip

\textbf{Question 4}

\medskip

On considère la fonction $f$ définie pour tout réel $x$ par $f(x)=\sin x -\cos x$. Parmi les quatre propositions suivantes, une seule est correcte. Laquelle ?

\medskip

\begin{tabularx}{\linewidth}{*{1}{X}}
\textbf{a.~~} $f$ est une fonction paire. \\\textbf{b.~~}$f$ est une fonction impaire. \\\textbf{c.~~}$f$ n'est ni paire ni impaire.\\ \textbf{d.~~}$f (0)=0$ .
\end{tabularx}

\medskip

\textbf{Question 5}

\medskip

Dans le plan muni d'un repère, on considère la droite $(d)$ d'équation : $5x-2y+8=0$.  

La droite $(d)$ a pour coefficient directeur :

\medskip

\begin{tabularx}{\linewidth}{*{4}{X}}
\textbf{a.~~} $\vv{u}(2 ;5)$ &\textbf{b.~~}$\dfrac{5}{2}$ &\textbf{c.~~}$\dfrac{ 2}{5}$& \textbf{d.~~} $ -2$.
\end{tabularx}



\vspace{0.75cm}

\end{document}
